\setcounter{page}{2}

\chapter*{РЕФЕРАТ}
Дипломная работа, 53 с., 8 рис., 15 источников.

\textbf{Ключевые слова}: РАСПРЕДЕЛЕННЫЕ СИСТЕМЫ; АЛГОРИТМ КОНСЕНСУСА; RAFT; ОТКАЗОУСТОЙЧИВОСТЬ; ДИНАМИЧЕСКАЯ РЕКОНФИГУРАЦИЯ; ХРАНИЛИЩЕ КЛЮЧ-ЗНАЧЕНИЕ; ЛИНЕАРИЗУЕМОСТЬ.

\textbf{Объект исследования}: распределенные системы хранения и обработки данных в условиях ненадежной сетевой среды.

\textbf{Цель работы}: проектирование архитектуры и программная реализация отказоустойчивого распределенного хранилища «ключ-значение» на основе алгоритма Raft, поддерживающего строгую согласованность данных и динамическое изменение топологии кластера без прерывания обслуживания.

\textbf{Методы исследования}: системный анализ, теория распределенных вычислений, математическое моделирование, объектно-ориентированное программирование, методы хаос-инжиниринга и нагрузочного тестирования.

\textbf{Результаты}: спроектирована и реализована многоуровневая архитектура распределенного хранилища. Разработан механизм строгой линеаризуемости чтений, исключающий возврат устаревших данных при сетевых разделениях. Реализован алгоритм безопасного динамического добавления и удаления узлов с использованием промежуточной фазы «наблюдателя». Разработан внутрипроцессный эмулятор сети, с помощью которого проведено тестирование системы методами хаос-инжиниринга, подтвердившее отказоустойчивость системы и стопроцентную защиту от потери или дублирования транзакций при аппаратных сбоях.

\textbf{Области применения}: облачные вычисления, микросервисная архитектура, системы обнаружения сервисов, высоконадежные реестры конфигураций.

\chapter*{РЭФЕРАТ}
Дыпломная праца, 53 с., 8 мал., 15 крынiц.

\textbf{Ключавыя словы}: РАЗМЕРКАВАНЫЯ СІСТЭМЫ; АЛГАРЫТМ КАНСЭНСУСУ; RAFT; АДМОВАЎСТОЙЛІВАСЦЬ; ДЫНАМІЧНАЯ РЭКАНФІГУРАЦЫЯ; СХОВІШЧА КЛЮЧ-ЗНАЧЭННЕ; ЛІНЕАРЫЗУ\-EМАСЦЬ.

\textbf{Аб’ект даследавання}: размеркаваныя сістэмы захоўвання і апрацоўкі дадзеных ва ўмовах ненадзейнага сеткавага асяроддзя.

\textbf{Мэта працы}: праектаванне архітэктуры і праграмная рэалізацыя адмоваўстойлівага размеркаванага сховішча «ключ-значэнне» на аснове алгарытму Raft, якое падтрымлівае строгую ўзгодненасць дадзеных і дынамічнае змяненне тапалогіі кластара без прыпынку абслугоўвання.

\textbf{Метады даследавання}: сістэмны аналіз, тэорыя размеркаваных вылічэнняў, матэматычнае мадэляванне, аб'ектна-арыентаванае праграмаванне, метады хаос-інжынірынгу і нагрузачнага тэсціравання.

\textbf{Вынікі}: спраектавана і рэалізавана шматузроўневая архітэктура размеркаванага сховішча. Распрацаваны механізм строгай лінеарызуемасці чытанняў, які выключае вяртанне састарэлых дадзеных пры сеткавых раздзяленнях. Рэалізаваны алгарытм бяспечнага дынамічнага дадання і выдалення вузлоў з выкарыстаннем прамежкавай фазы «назіральніка». Распрацаваны ўнутрыпрацэсны эмулятар сеткі, з дапамогай якога праведзена тэсціраванне сістэмы метадамі хаос-інжынірынгу, што пацвердзіла адмоваўстойлівасць сістэмы і стопрацэнтную абарону ад страты або дублявання транзакцый пры апаратных збоях.

\textbf{Сферы выкарыстання}: воблачныя вылічэнні, мікрасэрвісная архітэктура, сістэмы выяўлення сэрвісаў, высоканадзейныя рэестры канфігурацый.

\chapter*{ABSTRACT}
Diploma thesis, 53 p., 8 fig., 15 sources.

\textbf{Keywords}: DISTRIBUTED SYSTEMS; CONSENSUS ALGORITHM; RAFT; FAULT TOLERANCE; DYNAMIC RECONFIGURATION; KEY-VALUE STORE; LINEARIZABILITY.

\textbf{Object of research}: distributed data storage and processing systems in an unreliable network environment.

\textbf{Aim of the work}: architectural design and software implementation of a fault-tolerant distributed Key-Value store based on the Raft algorithm, supporting strict data consistency and dynamic cluster topology changes without service downtime.

\textbf{Research methods}: systematic analysis, theory of distributed computing, mathe\-matical modeling, object-oriented programming, chaos engineering, and load testing methods.

\textbf{Results}: a multi-layer architecture of a distributed store was designed and implemented. A strict read linearizability mechanism was developed, preventing the return of stale data during network partitions. A safe algorithm for dynamically adding and removing nodes was implemented using an intermediate "observer" phase. An in-process network emulator was developed, through which the system was tested using chaos engineering methods, confirming the fault tolerance of the system and one hundred percent protection against transaction loss or duplication during hardware failures.

\textbf{Areas of application}: cloud computing, microservice architecture, service discovery systems, highly reliable configuration registries.
